\section*{Capítulo \ref{ch:g12:buffers-predominance} --- Soluções-tampão e diagramas de predominância}

\begin{solution}{exo:g12:buffers-predominance:1}
\omterm{def:g9:acids-bases-ph:ph}{pH} 2: \ce{CH3COOH}. \omterm{def:g9:acids-bases-ph:ph}{pH} 4.76: os dois em quantidades iguais. \omterm{def:g9:acids-bases-ph:ph}{pH} 7: \ce{CH3COO-}.
\end{solution}

\begin{solution}{exo:g12:buffers-predominance:2}
\omterm{def:g9:acids-bases-ph:ph}{pH} 4: $1/(1 + 10^{0.76}) = 0.15$. \omterm{def:g9:acids-bases-ph:ph}{pH} 6: $1/(1 + 10^{-1.24}) = 0.95$.
\end{solution}

\begin{solution}{exo:g12:buffers-predominance:3}
Amarelo; verde (dentro da sua faixa); azul.
\end{solution}

\begin{solution}{exo:g12:buffers-predominance:4}
A \omterm{def:g9:acids-bases-ph:ph}{pH} 7.4, abaixo de 9.26: \ce{NH4+}. A \omterm{def:g9:acids-bases-ph:ph}{pH} 11: \ce{NH3}.
\end{solution}

\begin{solution}{exo:g12:buffers-predominance:5}
Um indicador é ele próprio um \omterm{def:g12:ka-and-pka:bronsted}{par ácido--base}: em grande quantidade, ele
mudaria o \omterm{def:g9:acids-bases-ph:ph}{pH} que deve mostrar.
\end{solution}

\begin{solution}{exo:g12:buffers-predominance:6}
$10^{\mathrm{pH} - 4.76}$: 0.1; 1; 10.
\end{solution}

\begin{solution}{exo:g12:buffers-predominance:7}
Fenolftaleína (8.0--10.0) ou azul de timol (8.0--9.1): as suas faixas
contêm 8.7. O vermelho de metila vira entre 4.2 e 6.3, longe de 8.7: ele
já teria mudado muito antes.
\end{solution}

\begin{solution}{exo:g12:buffers-predominance:8}
$\mathrm{pH} = 4.76 + \log(0.10/0.20) = 4.76 - 0.30 = 4.46$.
\end{solution}

\begin{solution}{exo:g12:buffers-predominance:9}
\omterm{def:g9:acids-bases-ph:ph}{pH} 1: o \omterm{def:g9:ions:ion}{cátion} \ce{H3N+-CH2-COOH}, carga $+1$. \omterm{def:g9:acids-bases-ph:ph}{pH} 6: o zwitteríon, carga
total 0. \omterm{def:g9:acids-bases-ph:ph}{pH} 12: o \omterm{def:g9:ions:ion}{ânion} \ce{H2N-CH2-COO-}, carga $-1$.
\end{solution}

\begin{solution}{exo:g12:buffers-predominance:10}
Água: de 7.0 para $-\log(\num{1.0e-4}/0.101) = 3.0$, uma queda de 4
unidades. Tampão: de 4.76 para $4.76 + \log(9.9/10.1) = 4.75$, uma queda
de 0.01.
\end{solution}

\begin{solution}{exo:g12:buffers-predominance:11}
\ce{NH4+}/\ce{NH3} ($pK_a = 9.26$). Razão
$[\ce{NH3}]/[\ce{NH4+}] = 10^{9.5 - 9.26} = 1.7$.
\end{solution}

\begin{solution}{exo:g12:buffers-predominance:12}
O \omterm{def:g9:acids-bases-ph:ph}{pH} cai uma unidade quando a razão chega a $1/10$:
$(10 - n)/(10 + n) = 0.1$, logo $n = \qty{8.2}{mmol}$. Um tampão que
contém mais de cada forma pode absorver mais ácido ou base para a mesma
mudança da razão: ele tem uma capacidade maior.
\end{solution}

\begin{solution}{exo:g12:buffers-predominance:13}
Razão $10^{5.0 - 4.76} = 1.74$; o ácido soma \qty{10}{mmol}, logo
\qty{17.4}{mmol} de etanoato, isto é, \qty{174}{mL} da solução.
\end{solution}

\begin{solution}{exo:g12:buffers-predominance:14}
Ele tem dois \omterm{def:g12:ka-and-pka:bronsted}{pares ácido--base}, com dois $pK_a$, daí duas mudanças de
cor: vermelho a \omterm{def:g9:acids-bases-ph:ph}{pH} 1, amarelo a \omterm{def:g9:acids-bases-ph:ph}{pH} 5, azul a \omterm{def:g9:acids-bases-ph:ph}{pH} 10.
\end{solution}

\begin{solution}{exo:g12:buffers-predominance:15}
A \omterm{def:g10:concentration-and-dilution:dilution}{diluição} divide as duas concentrações por 10: a razão e, portanto, o
\omterm{def:g9:acids-bases-ph:ph}{pH} (4.76) não mudam. Um \omterm{def:g12:ka-and-pka:strong-acid}{ácido forte} diluído dez vezes sobe uma unidade
de \omterm{def:g9:acids-bases-ph:ph}{pH}.
\end{solution}

\begin{solution}{pb:g12:buffers-predominance:1}
\textbf{1.} \ce{CO2 + 2H2O <=> HCO3- + H3O+}.

\textbf{2.} O dióxido de carbono dissolvido (com a água) é o ácido, o
hidrogenocarbonato é a base.

\textbf{3.} O valor da tabela é para \qty{25}{\celsius} em água pura; o
sangue está a \qty{37}{\celsius} e cheio de outros \omterm{def:g9:ions:ion}{íons}, que mudam a
constante (e os fisiologistas contam todo o dióxido de carbono
dissolvido).

\textbf{4.} A \qty{24}{mmol/L}, $0.024 \times 5.0 = \qty{0.12}{mol}$.

\textbf{5.} Um eixo de \omterm{def:g9:acids-bases-ph:ph}{pH} cortado em 6.1: \ce{CO2} abaixo, \ce{HCO3-}
acima.

\textbf{6.} A 7.4, acima de 6.1: o hidrogenocarbonato.

\textbf{7.} Levemente básico: o seu \omterm{def:g9:acids-bases-ph:ph}{pH} está acima de 7.

\textbf{8.} $10^{7.4 - 6.1} = 10^{1.3} = 20$.

\textbf{9.} $24 / 20 = \qty{1.2}{mmol/L}$.

\textbf{10.} $10^{1.28} = 19$ e $10^{1.32} = 21$.

\textbf{11.} $20/21 = \qty{95}{\percent}$.

\textbf{12.} $10^{7.6 - 6.1} = 10^{1.5} = 32$.

\textbf{13.} O \omterm{def:g9:ions:ion}{íon} hidrogenocarbonato:
\ce{HCO3- + H3O+ -> CO2 + 2H2O}.

\textbf{14.} $10^{7.1 - 6.1} = 10$.

\textbf{15.} De $24/20 = 1.2$ para $24/10 = \qty{2.4}{mmol/L}$: dobra. Os
pulmões respiram mais depressa e mais fundo, para eliminar o excesso de
dióxido de carbono.

\textbf{16.} Um \omterm{def:g12:ka-and-pka:strong-acid}{ácido fraco} e a sua base podem cada um absorver o que é
acrescentado, e o \omterm{def:g9:acids-bases-ph:ph}{pH} só depende da sua razão; um \omterm{def:g12:ka-and-pka:strong-acid}{ácido forte} só
empurraria o \omterm{def:g9:acids-bases-ph:ph}{pH} num sentido.

\textbf{17.} Nada: a razão e, portanto, o \omterm{def:g9:acids-bases-ph:ph}{pH} continuam os mesmos.

\textbf{18.} Cerca de 20.
\end{solution}
